\documentclass[aspectratio=169]{beamer}

% Tema minimalista
\usetheme{default}
\usecolortheme{whale} % tonos azules

% Opciones para notas del expositor (descomentar para ver)
% \setbeameroption{show notes on second screen=right}

% Colores personalizados
\definecolor{AzulOscuro}{RGB}{10, 30, 60}
\definecolor{AzulMedio}{RGB}{0, 102, 204}
\definecolor{GrisClaro}{RGB}{240, 240, 240}
\definecolor{GrisTexto}{RGB}{60, 60, 60}
\definecolor{Negro}{RGB}{0, 0, 0}

% Configuración de colores del tema
\setbeamercolor{palette primary}{bg=AzulOscuro,fg=white}
\setbeamercolor{palette secondary}{bg=AzulMedio,fg=white}
\setbeamercolor{palette tertiary}{bg=AzulOscuro,fg=white}
\setbeamercolor{palette quaternary}{bg=AzulOscuro,fg=white}
\setbeamercolor{structure}{fg=AzulMedio}
\setbeamercolor{normal text}{fg=GrisTexto}
\setbeamercolor{background canvas}{bg=white}
\setbeamercolor{frametitle}{bg=white,fg=AzulOscuro}
\setbeamercolor{title}{bg=white,fg=AzulOscuro}

% Configuración de fuentes y estilo
\usefonttheme{professionalfonts}
\usefonttheme[onlymath]{serif}
\setbeamerfont{title}{size=\huge, series=\bfseries}
\setbeamerfont{frametitle}{size=\Large, series=\bfseries}

% Limpiar plantilla para que sea más minimalista
\setbeamertemplate{navigation symbols}{}
\setbeamertemplate{footline}[frame number]
\setbeamertemplate{itemize items}[circle]
\setbeamertemplate{frametitle}{
    \vspace*{0.5cm}
    \insertframetitle
    \vspace*{0.2cm}
}

% Paquetes adicionales
\usepackage[utf8]{inputenc}
\usepackage[spanish]{babel}
\usepackage{graphicx}
\usepackage{booktabs}
\usepackage{tikz}
\usepackage{array}

\title{Democra}
\subtitle{Plataforma SaaS Multi-Tenant para la gestión integral de organizaciones de voluntariado}
\author{Sustentación de Proyecto - Ingeniería de Sistemas}
\date{\today}
\institute{Arquitectura de Software}

\begin{document}

% DIAPOSITIVA 1: Portada
\begin{frame}[plain]
    \titlepage
    \begin{center}
        % Reemplazar con el logo real de la universidad
        % \includegraphics[height=1.5cm]{img/logo_universidad.png} 
    \end{center}
    \note{
        Objetivo: Presentar el proyecto y contextualizar la exposición.\\
        Tiempo: 30 segundos.
    }
\end{frame}

% DIAPOSITIVA 2: Agenda
\begin{frame}{Agenda}
    \begin{columns}[T]
        \begin{column}{0.5\textwidth}
            \begin{itemize}
                \item Problema
                \item Solución propuesta
                \item Metodología DDS
                \item Arquitectura General
                \item Stack Tecnológico
            \end{itemize}
        \end{column}
        \begin{column}{0.5\textwidth}
            \begin{itemize}
                \item Desarrollo
                \item Calidad del Software
                \item Demostración en Vivo
                \item Conclusiones
            \end{itemize}
        \end{column}
    \end{columns}
    \note{
        Objetivo: Mostrar el recorrido de la exposición.\\
        Tiempo: 30 segundos.
    }
\end{frame}

% DIAPOSITIVA 3: Problema
\begin{frame}{Problema}
    \begin{columns}[T]
        \begin{column}{0.4\textwidth}
            \vspace{1cm}
            \begin{itemize}
                \item Gestión descentralizada en ONG.
                \item Procesos manuales ineficientes.
                \item Falta de control de voluntarios.
                \item Datos fragmentados y vulnerables.
            \end{itemize}
        \end{column}
        \begin{column}{0.6\textwidth}
            \centering
            % \includegraphics[width=\textwidth,height=0.7\textheight,keepaspectratio]{img/problema.png}
            \begin{tikzpicture}
                \node[draw=GrisClaro, fill=GrisClaro, minimum width=6cm, minimum height=4cm, text=GrisTexto] {Imagen del Problema (main.pdf)};
            \end{tikzpicture}
        \end{column}
    \end{columns}
    \note{
        Objetivo: Explicar la necesidad que originó Democra.\\
        Tiempo: 1 minuto.
    }
\end{frame}

% DIAPOSITIVA 4: Solución propuesta
\begin{frame}{Solución propuesta}
    \begin{columns}[T]
        \begin{column}{0.5\textwidth}
            \vspace{1cm}
            \begin{itemize}
                \item Plataforma unificada en la nube.
                \item Arquitectura SaaS Multi-Tenant.
                \item Gestión integral automatizada.
                \item Alta disponibilidad y seguridad.
            \end{itemize}
        \end{column}
        \begin{column}{0.5\textwidth}
            \centering
            % \includegraphics[width=\textwidth,height=0.7\textheight,keepaspectratio]{img/arquitectura_funcional.png}
            \begin{tikzpicture}
                \node[draw=GrisClaro, fill=GrisClaro, minimum width=5cm, minimum height=4cm, text=GrisTexto, align=center] {Diagrama\\Arquitectura Funcional};
            \end{tikzpicture}
        \end{column}
    \end{columns}
    \note{
        Objetivo: Explicar qué hace Democra.\\
        Tiempo: 1 minuto.
    }
\end{frame}

% DIAPOSITIVA 5: ¿Qué es Democra?
\begin{frame}{¿Qué es Democra?}
    \begin{columns}[T]
        \begin{column}{0.3\textwidth}
            \vspace{1cm}
            \begin{itemize}
                \item Entorno multi-organización.
                \item Panel de control centralizado.
                \item Sistema escalable y robusto.
            \end{itemize}
        \end{column}
        \begin{column}{0.7\textwidth}
            \centering
            % \includegraphics[width=\textwidth,height=0.7\textheight,keepaspectratio]{img/sistema_funcionando.png}
            \begin{tikzpicture}
                \node[draw=GrisClaro, fill=GrisClaro, minimum width=7cm, minimum height=5cm, text=GrisTexto, align=center] {Captura del\\Sistema Funcionando};
            \end{tikzpicture}
        \end{column}
    \end{columns}
    \note{
        Objetivo: Mostrar el sistema desde una perspectiva general.\\
        Tiempo: 1 minuto.
    }
\end{frame}

% DIAPOSITIVA 6: Metodología DDS
\begin{frame}{Metodología DDS}
    \begin{columns}[T]
        \begin{column}{0.4\textwidth}
            \vspace{1cm}
            \begin{itemize}
                \item Diseño Guiado por el Dominio.
                \item Enfoque en reglas de negocio.
                \item Escalabilidad técnica asegurada.
                \item Mantenibilidad a largo plazo.
            \end{itemize}
        \end{column}
        \begin{column}{0.6\textwidth}
            \centering
            % \includegraphics[width=\textwidth,height=0.7\textheight,keepaspectratio]{img/dds.png}
            \begin{tikzpicture}
                \node[draw=GrisClaro, fill=GrisClaro, minimum width=6cm, minimum height=4cm, text=GrisTexto, align=center] {Imagen Metodología DDS};
            \end{tikzpicture}
        \end{column}
    \end{columns}
    \note{
        Objetivo: Presentar la metodología utilizada.\\
        Tiempo: 2 minutos.
    }
\end{frame}

% DIAPOSITIVA 7: Flujo DDS
\begin{frame}{Flujo DDS}
    \centering
    \vspace{0.5cm}
    \begin{tikzpicture}[node distance=1.6cm, auto,
        box/.style={rectangle, draw=AzulOscuro, fill=white, thick, minimum width=2.8cm, minimum height=0.8cm, align=center, font=\footnotesize\sffamily},
        arrow/.style={->, thick, AzulMedio}]
        
        \node[box] (prob) {Problema};
        \node[box, right of=prob, xshift=1.6cm] (inv) {Investigación};
        \node[box, right of=inv, xshift=1.6cm] (dis) {Diseño};
        \node[box, right of=dis, xshift=1.6cm] (arq) {Arquitectura};
        
        \node[box, below of=arq, yshift=-0.5cm] (imp) {Implementación};
        \node[box, left of=imp, xshift=-1.6cm] (test) {Testing};
        \node[box, left of=test, xshift=-1.6cm] (doc) {Documentación};
        \node[box, left of=doc, xshift=-1.6cm] (desp) {Despliegue};
        
        \draw[arrow] (prob) -- (inv);
        \draw[arrow] (inv) -- (dis);
        \draw[arrow] (dis) -- (arq);
        \draw[arrow] (arq) -- (imp);
        \draw[arrow] (imp) -- (test);
        \draw[arrow] (test) -- (doc);
        \draw[arrow] (doc) -- (desp);
    \end{tikzpicture}
    
    \vspace{1cm}
    \begin{itemize}
        \item Proceso iterativo e incremental.
        \item Calidad asegurada en cada fase.
    \end{itemize}
    \note{
        Objetivo: Explicar cada etapa del desarrollo.\\
        Tiempo: 2 minutos.
    }
\end{frame}

% DIAPOSITIVA 8: Arquitectura General
\begin{frame}{Arquitectura General}
    \begin{columns}[T]
        \begin{column}{0.3\textwidth}
            \vspace{1cm}
            \begin{itemize}
                \item Separación de capas.
                \item Desacoplamiento.
                \item Servicios modulares.
                \item Datos unidireccionales.
            \end{itemize}
        \end{column}
        \begin{column}{0.7\textwidth}
            \centering
            % \includegraphics[width=\textwidth,height=0.7\textheight,keepaspectratio]{img/arquitectura.png}
            \begin{tikzpicture}
                \node[draw=GrisClaro, fill=GrisClaro, minimum width=7cm, minimum height=5cm, text=GrisTexto, align=center] {Diagrama\\Arquitectura General};
            \end{tikzpicture}
        \end{column}
    \end{columns}
    \note{
        Objetivo: Mostrar la arquitectura completa.\\
        Tiempo: 2 minutos.
    }
\end{frame}

% DIAPOSITIVA 9: Arquitectura del Código
\begin{frame}{Arquitectura del Código}
    \begin{columns}[T]
        \begin{column}{0.5\textwidth}
            \vspace{1cm}
            \begin{itemize}
                \item Modularidad en estructura.
                \item Carpetas por dominio.
                \item Reusabilidad de componentes.
                \item Código mantenible.
            \end{itemize}
        \end{column}
        \begin{column}{0.5\textwidth}
            \centering
            \begin{tikzpicture}[node distance=1cm,
                box/.style={rectangle, draw=AzulOscuro, thick, minimum width=4cm, minimum height=0.6cm, align=center, font=\footnotesize},
                arrow/.style={->, thick, AzulMedio}]
                
                \node[box, fill=white] (user) {Usuario};
                \node[box, fill=GrisClaro, below of=user] (front) {Frontend};
                \node[box, fill=white, below of=front] (hooks) {Hooks};
                \node[box, fill=GrisClaro, below of=hooks] (serv) {Servicios};
                \node[box, fill=white, below of=serv] (api) {API};
                \node[box, fill=GrisClaro, below of=api] (supa) {Supabase (PostgreSQL)};
                
                \draw[arrow] (user) -- (front);
                \draw[arrow] (front) -- (hooks);
                \draw[arrow] (hooks) -- (serv);
                \draw[arrow] (serv) -- (api);
                \draw[arrow] (api) -- (supa);
            \end{tikzpicture}
        \end{column}
    \end{columns}
    \note{
        Objetivo: Explicar cómo se organizó el proyecto.\\
        Tiempo: 2 minutos.
    }
\end{frame}

% DIAPOSITIVA 10: Stack Tecnológico
\begin{frame}{Stack Tecnológico}
    \centering
    \renewcommand{\arraystretch}{1.6}
    \begin{tabular}{@{}ll@{}}
        \toprule
        \textbf{Categoría} & \textbf{Tecnologías} \\
        \midrule
        \textcolor{AzulMedio}{Frontend} & React, Vite, Tailwind CSS \\
        \textcolor{AzulMedio}{Backend} & Node.js, Express \\
        \textcolor{AzulMedio}{Base de Datos} & PostgreSQL (Supabase) \\
        \textcolor{AzulMedio}{Testing} & Vitest, Jest, Supertest \\
        \textcolor{AzulMedio}{Seguridad} & JWT, RLS, Helmet, OTP \\
        \textcolor{AzulMedio}{Documentación} & Swagger, Graphify \\
        \bottomrule
    \end{tabular}
    \note{
        Objetivo: Mostrar todas las tecnologías.\\
        Tiempo: 2 minutos.
    }
\end{frame}

% DIAPOSITIVA 11: Arquitectura de Datos
\begin{frame}{Arquitectura de Datos (Multi-Tenant)}
    \begin{columns}[T]
        \begin{column}{0.4\textwidth}
            \vspace{1cm}
            \begin{itemize}
                \item Aislamiento de datos.
                \item Seguridad por inquilino.
                \item Rendimiento optimizado.
                \item Escalabilidad horizontal.
            \end{itemize}
        \end{column}
        \begin{column}{0.6\textwidth}
            \centering
            % \includegraphics[width=\textwidth,height=0.7\textheight,keepaspectratio]{img/arquitectura_datos.png}
            \begin{tikzpicture}
                \node[draw=GrisClaro, fill=GrisClaro, minimum width=6cm, minimum height=4cm, text=GrisTexto, align=center] {Diagrama\\Base de Datos};
            \end{tikzpicture}
        \end{column}
    \end{columns}
    \note{
        Objetivo: Explicar Multi-Tenant.\\
        Tiempo: 2 minutos.
    }
\end{frame}

% DIAPOSITIVA 12: Seguridad
\begin{frame}{Seguridad}
    \begin{columns}[T]
        \begin{column}{0.5\textwidth}
            \vspace{1cm}
            \begin{itemize}
                \item \textbf{RBAC:} Control de accesos.
                \item \textbf{JWT:} Autenticación segura.
                \item \textbf{RLS:} Políticas a nivel de fila.
                \item \textbf{OTP:} Accesos un solo uso.
                \item \textbf{Helmet:} Protección HTTP.
            \end{itemize}
        \end{column}
        \begin{column}{0.5\textwidth}
            \centering
            % \includegraphics[width=\textwidth,height=0.6\textheight,keepaspectratio]{img/seguridad.png}
            \begin{tikzpicture}
                \node[draw=GrisClaro, fill=GrisClaro, minimum width=5cm, minimum height=4cm, text=GrisTexto, align=center] {Esquema de\\Seguridad};
            \end{tikzpicture}
        \end{column}
    \end{columns}
    \note{
        Objetivo: Mostrar cómo se protege el sistema.\\
        Tiempo: 1 minuto.
    }
\end{frame}

% DIAPOSITIVA 13: Wireframes
\begin{frame}{Wireframes}
    \begin{columns}[T]
        \begin{column}{0.3\textwidth}
            \vspace{1cm}
            \begin{itemize}
                \item Diseño centrado en usuario.
                \item Iteraciones ágiles.
                \item Validación de UX/UI.
            \end{itemize}
        \end{column}
        \begin{column}{0.7\textwidth}
            \centering
            % \includegraphics[width=\textwidth,height=0.7\textheight,keepaspectratio]{img/wireframes.png}
            \begin{tikzpicture}
                \node[draw=GrisClaro, fill=GrisClaro, minimum width=7cm, minimum height=5cm, text=GrisTexto, align=center] {Imágenes de\\Wireframes};
            \end{tikzpicture}
        \end{column}
    \end{columns}
    \note{
        Objetivo: Mostrar la evolución del diseño.\\
        Tiempo: 1 minuto.
    }
\end{frame}

% DIAPOSITIVA 14: Implementación
\begin{frame}{Implementación}
    \begin{columns}[T]
        \begin{column}{0.4\textwidth}
            \vspace{1cm}
            \begin{itemize}
                \item Entorno de desarrollo.
                \item Control de versiones estricto.
                \item Revisión de código (PR).
                \item Estándares de calidad.
            \end{itemize}
        \end{column}
        \begin{column}{0.6\textwidth}
            \centering
            % \includegraphics[width=\textwidth,height=0.7\textheight,keepaspectratio]{img/vscode_captura.png}
            \begin{tikzpicture}
                \node[draw=GrisClaro, fill=GrisClaro, minimum width=6cm, minimum height=4cm, text=GrisTexto, align=center] {Captura de\\VSCode};
            \end{tikzpicture}
        \end{column}
    \end{columns}
    \note{
        Objetivo: Explicar cómo se construyó el sistema.\\
        Tiempo: 1 minuto.
    }
\end{frame}

% DIAPOSITIVA 15: Calidad del Software
\begin{frame}{Calidad del Software}
    \begin{columns}[T]
        \begin{column}{0.5\textwidth}
            \vspace{0.5cm}
            \begin{itemize}
                \item Estándar \textbf{ISO 25010}.
                \item \textbf{Vitest:} Pruebas unitarias UI.
                \item \textbf{Jest:} Lógica del backend.
                \item \textbf{Supertest:} Integración API.
                \item Ejecución continua (CI).
            \end{itemize}
        \end{column}
        \begin{column}{0.5\textwidth}
            \centering
            \begin{tikzpicture}[node distance=1.5cm,
                box/.style={rectangle, draw=AzulOscuro, thick, minimum width=3.5cm, minimum height=0.8cm, align=center, font=\footnotesize},
                arrow/.style={->, thick, AzulMedio}]
                
                \node[box, fill=white] (unit) {Pruebas Unitarias};
                \node[box, fill=GrisClaro, below of=unit] (int) {Pruebas Integración};
                \node[box, fill=white, below of=int] (cov) {Métricas de Cobertura};
                
                \draw[arrow] (unit) -- (int);
                \draw[arrow] (int) -- (cov);
            \end{tikzpicture}
        \end{column}
    \end{columns}
    \note{
        Objetivo: Mostrar la estrategia de calidad.\\
        Tiempo: 2 minutos.
    }
\end{frame}

% DIAPOSITIVA 16: Cobertura y Testing
\begin{frame}{Cobertura y Testing}
    \begin{columns}[T]
        \begin{column}{0.5\textwidth}
            \centering
            \textbf{Resultados de Pruebas}\\
            \vspace{0.3cm}
            % \includegraphics[width=0.9\textwidth,height=0.5\textheight,keepaspectratio]{img/captura_tests.png}
            \begin{tikzpicture}
                \node[draw=GrisClaro, fill=GrisClaro, minimum width=4cm, minimum height=3.5cm, text=GrisTexto, align=center] {Captura\\Tests};
            \end{tikzpicture}
        \end{column}
        \begin{column}{0.5\textwidth}
            \centering
            \textbf{Reporte de Cobertura}\\
            \vspace{0.3cm}
            % \includegraphics[width=0.9\textwidth,height=0.5\textheight,keepaspectratio]{img/captura_coverage.png}
            \begin{tikzpicture}
                \node[draw=GrisClaro, fill=GrisClaro, minimum width=4cm, minimum height=3.5cm, text=GrisTexto, align=center] {Captura\\Coverage};
            \end{tikzpicture}
        \end{column}
    \end{columns}
    \note{
        Objetivo: Explicar la cobertura y las pruebas.\\
        Tiempo: 2 minutos.
    }
\end{frame}

% DIAPOSITIVA 17: Documentación Técnica
\begin{frame}{Documentación Técnica}
    \begin{columns}[T]
        \begin{column}{0.4\textwidth}
            \vspace{1cm}
            \begin{itemize}
                \item \textbf{Swagger:} Contratos API.
                \item \textbf{Graphify:} Grafo de conocimiento.
                \item Trazabilidad completa.
                \item Integración continua.
            \end{itemize}
        \end{column}
        \begin{column}{0.6\textwidth}
            \centering
            % \includegraphics[width=\textwidth,height=0.7\textheight,keepaspectratio]{img/captura_swagger.png}
            \begin{tikzpicture}
                \node[draw=GrisClaro, fill=GrisClaro, minimum width=6cm, minimum height=4cm, text=GrisTexto, align=center] {Captura\\Swagger / Graphify};
            \end{tikzpicture}
        \end{column}
    \end{columns}
    \note{
        Objetivo: Mostrar la documentación.\\
        Tiempo: 1 minuto.
    }
\end{frame}

% DIAPOSITIVA 18: Demostración en Vivo
\begin{frame}{Demostración en Vivo}
    \centering
    \vspace{0.5cm}
    \Large \textcolor{AzulOscuro}{\textbf{Escenario de Demostración}}
    \normalsize
    \vspace{0.8cm}
    \begin{itemize}
        \item Mostrar estructura del proyecto (VSCode).
        \item Ejecutar suite de pruebas en terminal.
        \item Revisar reporte final de cobertura.
        \item Explorar contratos de API (Swagger).
        \item \textbf{Demostrar sistema Democra operando en vivo.}
    \end{itemize}
    \note{
        Objetivo: Mostrar el funcionamiento real.\\
        Tiempo: 5 minutos.
    }
\end{frame}

% DIAPOSITIVA 19: Resultados
\begin{frame}{Resultados}
    \centering
    \renewcommand{\arraystretch}{1.6}
    \begin{tabular}{@{}ll@{}}
        \toprule
        \textbf{Métrica} & \textbf{Resultado} \\
        \midrule
        \textcolor{AzulMedio}{Plataforma} & Operativa y Desplegada \\
        \textcolor{AzulMedio}{Arquitectura} & SaaS Multi-Tenant Implementada \\
        \textcolor{AzulMedio}{Seguridad} & RLS, RBAC y JWT Activos \\
        \textcolor{AzulMedio}{Cobertura} & Alta fiabilidad del código lograda \\
        \bottomrule
    \end{tabular}
    \note{
        Objetivo: Resumir los logros técnicos.\\
        Tiempo: 1 minuto.
    }
\end{frame}

% DIAPOSITIVA 20: Conclusiones
\begin{frame}{Conclusiones}
    \begin{columns}[T]
        \begin{column}{0.5\textwidth}
            \textcolor{AzulOscuro}{\textbf{Logros Alcanzados}}
            \vspace{0.3cm}
            \begin{itemize}
                \item Solución integral para ONG.
                \item Arquitectura escalable y segura.
                \item Código mantenible y testeado.
            \end{itemize}
        \end{column}
        \begin{column}{0.5\textwidth}
            \textcolor{AzulOscuro}{\textbf{Trabajo Futuro}}
            \vspace{0.3cm}
            \begin{itemize}
                \item Integración de Inteligencia Artificial.
                \item Módulos contables avanzados.
                \item Aplicación móvil nativa.
            \end{itemize}
        \end{column}
    \end{columns}
    \note{
        Objetivo: Cerrar la exposición.\\
        Tiempo: 1 minuto.
    }
\end{frame}

\end{document}
